\chapter{Codes}
\label{cha:app1_codes}
\begin{lstlisting}[language=python, caption = {General utility functions}, label = {lst:code-utility}, basicstyle=\ttfamily\scriptsize, % Shrinks the font size
    lineskip=-1pt,                   % Pulls the lines of code closer together
    aboveskip=10pt,                  % Reduces white space above the box
    belowskip=10pt                   % Reduces white space below the box
]
import numpy as np
from scipy.linalg import block_diag
from scipy.linalg import eigvalsh
from toqito.channel_ops import apply_channel
from toqito.matrices import gen_pauli_z as clock, gen_pauli_x as shift

np.set_printoptions(precision=5, suppress=True)
from fid_vectors import fid_vector_list

import jax
import jax.numpy as jnp

import numpy as np

def depol_jax(rho, p):
    d = rho.shape[0]
    return (1-p)* rho + p* (jnp.eye(d) / d) * jnp.trace(rho)

def compute_ptm_jax(channel_func_single_arg, d):
    basis_matrices = get_orthonormal_basis_first_identity(d) # Deine Funktion von vorhin
    basis_stack = jnp.array(basis_matrices)
    
    basis_flat = basis_stack.reshape(basis_stack.shape[0], -1)
    outputs_flat = outputs_stack.reshape(outputs_stack.shape[0], -1)
    
    ptm = jnp.conj(basis_flat) @ outputs_flat.T
    return ptm


def get_orthonormal_basis_first_identity(d: int):
    """
    Generates the d^2 Orthonormal Generalized Gell-Mann matrices.
    IMPORTANT: The normalized Identity is placed at index 0.
    
    Args:
        d (int): Dimension of the Hilbert space.
        
    Returns:
        list[np.ndarray]: List of basis matrices E_0, ..., E_{d^2-1}.
    """
    basis = []

    # E_0 = I / sqrt(d)
    basis.append(np.eye(d, dtype=complex) / np.sqrt(d))
    
    # 2. Symmetric Off-Diagonal
    for k in range(d):
        for j in range(k):
            mat = np.zeros((d, d), dtype=complex)
            mat[j, k] = 1
            mat[k, j] = 1
            basis.append(mat / np.sqrt(2))
            
    # 3. Anti-Symmetric Off-Diagonal
    for k in range(d):
        for j in range(k):
            mat = np.zeros((d, d), dtype=complex)
            mat[j, k] = -1j
            mat[k, j] = 1j
            basis.append(mat / np.sqrt(2))
            
    # 4. Diagonal (Traceless)
    for l in range(1, d):
        mat = np.zeros((d, d), dtype=complex)
        coeff = np.sqrt(1 / (l * (l + 1)))
        for m in range(l):
            mat[m, m] = coeff
        mat[l, l] = -l * coeff
        basis.append(mat)
        
    return basis

def compute_ptm(channel_func, d: int) -> np.ndarray:
    """
    Computes the Pauli Transfer Matrix (PTM) for a black-box channel.
    
    Args:
        channel_func (callable): A function T(rho) -> rho_prime.
        d (int): Dimension of the system.
        
    Returns:
        np.ndarray: The d^2 x d^2 matrix M where M_ij = Tr(E_i @ T(E_j)).
    """
    # 1. Get the basis E_0, ..., E_{d^2-1}
    basis = get_orthonormal_basis_first_identity(d)
    n = len(basis) # Should be d^2
    
    # 2. Initialize the empty PTM
    ptm = np.zeros((n, n), dtype=complex)
    
    # 3. Compute column by column (to minimize calls to channel_func)
    for j in range(n):
        E_j = basis[j]
        
        output_rho = channel_func(E_j)
        
        for i in range(n):
            E_i = basis[i]
            
            # M_ij = Tr(E_i @ T(E_j))
            # Note: Since basis is Hermitian, E_i^dagger = E_i.
            val = np.trace(E_i @ output_rho)
            ptm[i, j] = val
            
    return ptm

def compute_ptm_vectorized(channel_func, d: int) -> np.ndarray:
    """
    Computes the Pauli Transfer Matrix (PTM) using vectorization.
    """
    basis_flat = np.array([b.flatten() for b in basis_matrices])

    output_states = [channel_func(E_j) for E_j in basis_matrices]

    outputs_flat = np.array([rho.flatten() for rho in output_states])

    ptm = basis_flat.conj() @ outputs_flat.T
    
    return ptm

    


def heisenberg_weyl_toqito(x: int, z: int, d: int) -> np.ndarray:
    """
    Implements the Heisenberg-Weyl operator using toqito's generalized Pauli function.
    
    Toqito's `gen_pauli(k, 1, d)` corresponds to X.
    Toqito's `gen_pauli(k, 2, d)` corresponds to Z.
    
    However, toqito does not have a single function taking both x and z directly 
    as a combined W operator in one call, so we construct it similarly.
    """
    try:
        from toqito.matrices import gen_pauli
    except ImportError:
        return "Toqito not installed. Run `pip install toqito`."
    
    X_op = gen_pauli(x, 1, d)
    
    Z_op = gen_pauli(z, 2, d)
    
    return X_op @ Z_op


def depol_channel(rho):
    return apply_channel(rho,choi_depolarizing)

def top_two_abs_values(arr: np.ndarray, decimals: int = 5):
    """
    Analyzes a 1D numpy array to find the counts of the highest 
    and second-highest absolute values, robust to floating point errors.

    Args:
        arr (np.ndarray): Input array of shape (n,).
        decimals (int): Number of decimal places to round to. 
                        Default is 10, which handles most numerical noise.

    Returns:
        dict: Statistics about the top two absolute values.
    """
    # 1. Compute absolute values
    abs_vals = np.abs(arr)
    
    # 2. CRITICAL STEP: Round to handle floating point noise
    # This groups 0.300000001 and 0.299999999 together as 0.3
    abs_vals = np.round(abs_vals, decimals=decimals)
    
    # 3. Find unique values and their counts
    unique_vals, counts = np.unique(abs_vals, return_counts=True)
    
    # 4. Handle Edge Case: Empty Array
    if len(unique_vals) == 0:
        return {
            "first_max_val": None, "first_max_count": 0,
            "second_max_val": None, "second_max_count": 0
        }

    # 5. Extract First Highest (Last element in sorted array)
    first_max_val = unique_vals[-1]
    first_max_count = counts[-1]

    # 6. Extract Second Highest (Second to last element)
    if len(unique_vals) >= 2:
        second_max_val = unique_vals[-2]
        second_max_count = counts[-2]
    else:
        # Edge Case: All numbers have the same absolute value
        second_max_val = None
        second_max_count = 0

    return {
        "first_max_val": first_max_val,
        "first_max_count": first_max_count,
        "second_max_val": second_max_val,
        "second_max_count": second_max_count
    }
\end{lstlisting}

\begin{lstlisting}[language = python, caption = Implementation of the SIC-POVM condition, label = {lst:code-sic}, basicstyle=\ttfamily\scriptsize, % Shrinks the font size
    lineskip=-1pt,                   % Pulls the lines of code closer together
    aboveskip=10pt,                  % Reduces white space above the box
    belowskip=10pt                   % Reduces white space below the box
]

def get_sic_povm(d, fiducial_vector):
    """
    Generates a SIC-POVM from a fiducial vector using Weyl-Heisenberg operators.
    """
    X = shift(d)
    Z = clock(d)
    
    povm_elements = []
    # Displacement operators D_{a,b} = X^a Z^b (ignoring phase factors for POVM)
    for a in range(d):
        for b in range(d):
            # D = (X^a) @ (Z^b)
            op = np.linalg.matrix_power(X, a) @ np.linalg.matrix_power(Z, b)
            vec = op @ fiducial_vector
            # POVM element E_k = 1/d * |psi><psi|
            povm_elements.append((1/d) * np.outer(vec, vec.conj()))
            
    return povm_elements

def sic_condition_jax(channel_func, d):
    """
    channel_func: Deine 'channel_ready' Funktion (lambda), die nur ein Argument (rho) nimmt.
    """

    fid_vector = jnp.array(fid_vector_list[d-2]) 

    fid_vector = fid_vector / jnp.linalg.norm(fid_vector)
    
    sic_povm_np = get_sic_povm(d, fid_vector) 
    sic_povm = [jnp.array(M) for M in sic_povm_np]
    
    temp_sic = True

    for i in sic_povm:
        channel_out = channel_func(i)
        
        temp = (d + 1) * channel_out - jnp.eye(d) / d

        if not is_positive_semidefinite_jax(temp):
            temp_sic = False
            # Optional: break, um Zeit zu sparen, wenn einer fehlschlägt
            break 

    if temp_sic:
        print('SICs say channel is entanglement breaking')
    else:
        print('SICs can not make a statement')
        
    return temp_sic
\end{lstlisting}

\begin{lstlisting}[language = python, caption = {Optimization of one PPT channel with an SDP program using the Breuer-Hall map}, label = {lst:breuer}, basicstyle=\ttfamily\scriptsize, % Shrinks the font size
    lineskip=-1pt,                   % Pulls the lines of code closer together
    aboveskip=10pt,                  % Reduces white space above the box
    belowskip=10pt                   % Reduces white space below the box
]
def optimize_channel(target, J_fixed, rho_fixed, d, U):
    """
    Optimizes either T1 or T2 to minimize Tr(rho * M(T1, T2)).
    
    Parameters:
    - target: String, either 'T1' or 'T2' (the variable to optimize).
    - J_fixed: numpy array (d^2 x d^2), the Choi matrix of the fixed channel.
    - rho_fixed: numpy array (d^2 x d^2), the fixed state (eigenvector).
    - d: Integer, dimension of the subsystem (must be even, >= 4).
    
    Returns:
    - min_val: The minimum expectation value found.
    - J_opt: The Choi matrix of the optimized channel.
    """
    d2 = d**2
    P_reshuffle, P_PT, P_TrB = build_linear_maps(d)
    P_reshuffle_inv = P_reshuffle.T # Inverse of permutation is its transpose
    
    # Define the CVXPY Variable for the target Choi matrix
    J_var = cp.Variable((d2, d2), hermitian=True)
    
    # Convert fixed Choi to Superoperator
    K_fixed = (P_reshuffle @ J_fixed.flatten()).reshape((d2, d2), order = 'C')
    
    # Convert variable Choi to Superoperator
    K_var = cp.reshape(P_reshuffle @ cp.vec(J_var, order='C'), (d2, d2), order = 'C')
    
    # Compose channels (Matrix multiplication of superoperators)
    if target == 'T1':
        # T1 is var, T2 is fixed. K_composed = K_var @ K_fixed
        K_composed = K_var @ K_fixed
    elif target == 'T2':
        # T2 is var, T1 is fixed. K_composed = K_fixed @ K_var
        K_composed = K_fixed @ K_var
    else:
        raise ValueError("Target must be 'T1' or 'T2'")
        
    # Convert composed Superoperator back to Choi matrix
    J_composed = cp.reshape(P_reshuffle_inv @ cp.vec(K_composed, order='C'), (d2, d2))
    
    # --- Apply Breuer-Hall Map to J_composed ---
    #U = get_breuer_hall_U(d)
    I_U = np.kron(np.eye(d), U)
    
    # J_composed is TP, so Tr_B(J) x I = I_d^2
    term1 = np.eye(d2)
    term2 = J_composed
    
    # Apply (I x U) * J_composed^{T_B} * (I x U^dagger)
    J_comp_PT = cp.reshape(P_PT @ cp.vec(J_composed, order='C'), (d2, d2), order = 'C')
    term3 = I_U @ J_comp_PT @ I_U.conj().T
    
    M = term1 - term2 - term3
    
    # --- Objective ---
    objective = cp.Minimize(cp.real(cp.trace(rho_fixed @ M)))
    
    # --- Constraints for J_var to be a valid PPT Channel ---
    constraints = []
    constraints.append(J_var >> 0) # Completely Positive
    
    # Partial Transpose Constraint (PPT)
    J_var_PT = cp.reshape(P_PT @ cp.vec(J_var, order='C'), (d2, d2), order = 'C')
    constraints.append(J_var_PT >> 0)
    
    # Trace Preserving Constraint: Tr_B(J) = I_d
    TrB_J_var = cp.reshape(P_TrB @ cp.vec(J_var, order='C'), (d, d), order = 'C')
    constraints.append(TrB_J_var == np.eye(d))
    
    # --- Solve ---
    prob = cp.Problem(objective, constraints)
    prob.solve(solver=cp.SCS) # SCS or MOSEK recommended for large complex SDPs
    
    return prob.value, J_var.value

\end{lstlisting}

\begin{lstlisting}[language = python, caption = Testing the PPT$^2$ conjecture with SDPs, label={lst:testing-ppt-2}, basicstyle=\ttfamily\scriptsize, % Shrinks the font size
    lineskip=-1pt,                   % Pulls the lines of code closer together
    aboveskip=10pt,                  % Reduces white space above the box
    belowskip=10pt                   % Reduces white space below the box
]
    import numpy as np
import cvxpy as cp
from scipy.optimize import minimize
import autograd.numpy as anp
from autograd import value_and_grad
from scipy.linalg import qr
import scipy.sparse as sp

def get_random_breuer_hall_U(d):
    """
    Generates a random antisymmetric unitary matrix U for the Breuer-Hall map.
    """
    # 1. Construct the trivial antisymmetric unitary (your original Omega)
    Omega = np.zeros((d, d), dtype=complex)
    for i in range(d // 2):
        Omega[2*i, 2*i+1] = 1.0
        Omega[2*i+1, 2*i] = -1.0
        
    # 2. Generate a random complex Gaussian matrix
    G = np.random.randn(d, d) + 1j * np.random.randn(d, d)
    
    # 3. Extract a random unitary matrix V using QR decomposition
    V, _ = qr(G)
    
    # 4. Construct the random antisymmetric unitary
    # CRITICAL: Use V.T (transpose), NOT V.conj().T (Hermitian conjugate)
    U_random = V @ Omega @ V.T
    
    return U_random

def generate_random_density_matrix(dim):
    """
    Generiert eine zufällige Dichtematrix der Dimension (dim x dim).
    """
    G = np.random.randn(dim, dim) + 1j * np.random.randn(dim, dim)
    rho_unnormalized = G @ G.conj().T
    rho = rho_unnormalized / np.trace(rho_unnormalized)
    return rho

def generate_random_boundary_ppt_channel(d, P_PT, P_TrB):
    d2 = d**2

    G = np.random.randn(d2, d2) + 1j * np.random.randn(d2, d2)
    H = G + G.conj().T
    

    C_PT = cp.reshape(P_PT @ cp.vec(C, order='C'), (d2, d2), order='C')
    constraints.append(C_PT >> 0)
    
    TrB_C = cp.reshape(P_TrB @ cp.vec(C, order='C'), (d, d), order='C')
    constraints.append(TrB_C == np.eye(d))
    
    objective = cp.Minimize(cp.real(cp.trace(H @ C)))
    prob = cp.Problem(objective, constraints)
    prob.solve(solver=cp.SCS) 
    
    return C.value


def optimize_U_scalable(rho, J_PT, d, Omega):
    """
    Optimizes the antisymmetric unitary U for any dimension d using the Cayley Transform.
    """
    # Force all constants into autograd's tracked array format to prevent primitive leaks
    I_d = anp.eye(d, dtype=complex)
    Omega_anp = anp.array(Omega, dtype=complex)
    rho_anp = anp.array(rho, dtype=complex)
    J_PT_anp = anp.array(J_PT, dtype=complex)
    
    def objective(params):
        # 1. Map flat array to complex matrix Z
        Z_real = anp.reshape(params[:d**2], (d, d))
        Z_imag = anp.reshape(params[d**2:], (d, d))
        Z = Z_real + 1j * Z_imag
        
        # 2. Map Z to Hermitian H
        H = Z + anp.conj(anp.transpose(Z))
        
        # 3. Map Hermitian H to Unitary V via Cayley Transform (Bypassing expm)
        iH = 1j * H
        V = anp.matmul(I_d - iH, anp.linalg.inv(I_d + iH))
        
        # 4. Map V to Antisymmetric Unitary U
        U = anp.matmul(V, anp.matmul(Omega_anp, anp.transpose(V)))
        
        # 5. Construct tensor operator (autograd safe)
        I_U = anp.reshape(anp.einsum('ij,kl->ikjl', I_d, U), (d**2, d**2))
        I_U_dag = anp.conj(anp.transpose(I_U))
        
        # 6. Calculate objective trace
        val = anp.real(anp.trace(anp.matmul(rho_anp, anp.matmul(I_U, anp.matmul(J_PT_anp, I_U_dag)))))
        return -val

    x0 = np.random.randn(2 * d**2)
    
    res = minimize(value_and_grad(objective), x0, jac=True, method='BFGS')
    
    # Reconstruct optimal U purely in NumPy
    Z_real_opt = res.x[:d**2].reshape((d, d))
    Z_imag_opt = res.x[d**2:].reshape((d, d))
    Z_opt = Z_real_opt + 1j * Z_imag_opt
    H_opt = Z_opt + Z_opt.conj().T
    
    # Cayley Transform for the final output
    I_np = np.eye(d, dtype=complex)
    iH_opt = 1j * H_opt
    V_opt = (I_np - iH_opt) @ np.linalg.inv(I_np + iH_opt)
    U_opt = V_opt @ Omega @ V_opt.T
    
    return U_opt

def update_rho(M_matrix):
    """
    Updates the variational state rho by finding the eigenvector 
    corresponding to the smallest eigenvalue of M.
    
    Parameters:
    - M_matrix: A (d^2 x d^2) numpy array representing the Hermitian matrix M.
    
    Returns:
    - rho_new: The updated (d^2 x d^2) density matrix.
    - min_eval: The smallest eigenvalue (for tracking progress).
    """
    # 1. Calculate eigenvalues (w) and eigenvectors (v)
    # eigh guarantees w is sorted ascending, and v[:, i] is the i-th eigenvector
    w, v = np.linalg.eigh(M_matrix)
    
    # 2. Extract the smallest eigenvalue and its eigenvector
    min_eval = w[0]
    psi_min = v[:, 0]  # The column vector corresponding to w[0]
    
    # 3. Construct the new density matrix: rho = |psi_min><psi_min|
    # Using outer product of the vector and its complex conjugate
    rho_new = np.outer(psi_min, psi_min.conj())
    
    return rho_new, min_eval


def build_linear_maps(d):
    """
    Constructs the permutation matrices using ultra-fast, memory-safe 
    vectorized operations in COO format. No nested loops.
    """
    d2 = d**2
    d4 = d**4
    
    # -------------------------------------------------------------
    # 1. P_reshuffle (K_{kl, ij} = J_{ik, jl})
    # -------------------------------------------------------------
    # Generate all combinations of (i, j, k, l)
    i, j, k, l = np.mgrid[0:d, 0:d, 0:d, 0:d]
    
    # Flatten them to 1D arrays
    i = i.flatten()
    j = j.flatten()
    k = k.flatten()
    l = l.flatten()
    
    # Calculate row and column indices
    row_K = (k*d + l)*d2 + (i*d + j)
    col_J = (i*d + k)*d2 + (j*d + l)
    vals_K = np.ones(d4)
    
    # Build sparse matrix directly
    P_reshuffle = sp.coo_matrix((vals_K, (row_K, col_J)), shape=(d4, d4)).tocsr()
    
    # -------------------------------------------------------------
    # 2. P_PT (Partial Transpose: J^{T_B}_{ik, jl} = J_{il, jk})
    # -------------------------------------------------------------
    # The indices i, j, k, l are already defined above
    row_PT = (i*d + k)*d2 + (j*d + l)
    col_PT = (i*d + l)*d2 + (j*d + k)
    vals_PT = np.ones(d4)
    
    P_PT = sp.coo_matrix((vals_PT, (row_PT, col_PT)), shape=(d4, d4)).tocsr()
    
    # -------------------------------------------------------------
    # 3. P_TrB (Partial Trace on B: Tr_B(J)_{i,j} = sum_k J_{ik, jk})
    # -------------------------------------------------------------
    # For Partial Trace, we only iterate over i, j, k
    i_tr, j_tr, k_tr = np.mgrid[0:d, 0:d, 0:d]
    
    i_tr = i_tr.flatten()
    j_tr = j_tr.flatten()
    k_tr = k_tr.flatten()
    
    row_TrB = i_tr*d + j_tr
    col_TrB = (i_tr*d + k_tr)*d2 + (j_tr*d + k_tr)
    vals_TrB = np.ones(d**3) # Only d^3 elements
    
    P_TrB = sp.coo_matrix((vals_TrB, (row_TrB, col_TrB)), shape=(d2, d4)).tocsr()
    
    return P_reshuffle, P_PT, P_TrB

def get_breuer_hall_U(d):
    """Constructs the antisymmetric unitary U for the Breuer-Hall map."""
    U = np.zeros((d, d), dtype=complex)
    for i in range(d // 2):
        U[2*i, 2*i+1] = 1.0
        U[2*i+1, 2*i] = -1.0
    return U

def optimize_channel(target, J_fixed, rho_fixed, d, U):
    """
    Optimizes either T1 or T2 to minimize Tr(rho * M(T1, T2)).
    
    Parameters:
    - target: String, either 'T1' or 'T2' (the variable to optimize).
    - J_fixed: numpy array (d^2 x d^2), the Choi matrix of the fixed channel.
    - rho_fixed: numpy array (d^2 x d^2), the fixed state (eigenvector).
    - d: Integer, dimension of the subsystem (must be even, >= 4).
    
    Returns:
    - min_val: The minimum expectation value found.
    - J_opt: The Choi matrix of the optimized channel.
    """
    d2 = d**2
    P_reshuffle, P_PT, P_TrB = build_linear_maps(d)
    P_reshuffle_inv = P_reshuffle.T # Inverse of permutation is its transpose
    
    # Define the CVXPY Variable for the target Choi matrix
    J_var = cp.Variable((d2, d2), hermitian=True)
    
    # Convert fixed Choi to Superoperator
    K_fixed = (P_reshuffle @ J_fixed.flatten()).reshape((d2, d2), order = 'C')
    
    # Convert variable Choi to Superoperator
    K_var = cp.reshape(P_reshuffle @ cp.vec(J_var, order='C'), (d2, d2), order = 'C')
    
    # Compose channels (Matrix multiplication of superoperators)
    if target == 'T1':
        # T1 is var, T2 is fixed. K_composed = K_var @ K_fixed
        K_composed = K_var @ K_fixed
    elif target == 'T2':
        # T2 is var, T1 is fixed. K_composed = K_fixed @ K_var
        K_composed = K_fixed @ K_var
    else:
        raise ValueError("Target must be 'T1' or 'T2'")
        
    # Convert composed Superoperator back to Choi matrix
    J_composed = cp.reshape(P_reshuffle_inv @ cp.vec(K_composed, order='C'), (d2, d2), order = 'C')
    
    # --- Apply Breuer-Hall Map to J_composed ---
    #U = get_breuer_hall_U(d)
    I_U = np.kron(np.eye(d), U)
    
    # J_composed is TP, so Tr_B(J) x I = I_d^2
    term1 = np.eye(d2)
    term2 = J_composed
    
    # Apply (I x U) * J_composed^{T_B} * (I x U^dagger)
    J_comp_PT = cp.reshape(P_PT @ cp.vec(J_composed, order='C'), (d2, d2), order = 'C')
    term3 = I_U @ J_comp_PT @ I_U.conj().T
    
    M = term1 - term2 - term3
    
    # --- Objective ---
    objective = cp.Minimize(cp.real(cp.trace(rho_fixed @ M)))
    
    # --- Constraints for J_var to be a valid PPT Channel ---
    constraints = []
    constraints.append(J_var >> 0) # Completely Positive
    
    # Partial Transpose Constraint (PPT)
    J_var_PT = cp.reshape(P_PT @ cp.vec(J_var, order='C'), (d2, d2), order = 'C')
    constraints.append(J_var_PT >> 0)
    
    # Trace Preserving Constraint: Tr_B(J) = I_d
    TrB_J_var = cp.reshape(P_TrB @ cp.vec(J_var, order='C'), (d, d), order = 'C')
    constraints.append(TrB_J_var == np.eye(d))
    
    # --- Solve ---
    prob = cp.Problem(objective, constraints)
    prob.solve(solver=cp.SCS) # SCS or MOSEK recommended for large complex SDPs
    
    return prob.value, J_var.value

def calculate_M_matrix(J1, J2, U, d=6):
    d2 = d**2
    P_reshuffle, P_PT, P_TrB = build_linear_maps(d)
    
    # 1. Compose exactly like the SDP
    K1 = (P_reshuffle @ J1.flatten()).reshape((d2, d2), order='C')
    K2 = (P_reshuffle @ J2.flatten()).reshape((d2, d2), order='C')
    
    # K2 is applied first, then K1. Let's match your T1/T2 logic.
    # In your loop: min_val, J1 = optimize_channel('T1', J2...) -> K_composed = K_var @ K_fixed -> K1 @ K2
    # Wait, if T1 is applied after T2, it's K1 @ K2.
    K_composed = K1 @ K2 
    
    J_composed = (P_reshuffle.T @ K_composed.flatten()).reshape((d2, d2), order='C')
    
    # 2. PT exactly like the SDP
    J_comp_PT = (P_PT @ J_composed.flatten()).reshape((d2, d2), order='C')
    
    # 3. Breuer-Hall
    term1 = np.eye(d2)
    term2 = J_composed
    I_U = np.kron(np.eye(d), U)
    term3 = I_U @ J_comp_PT @ I_U.conj().T
    
    M = term1 - term2 - term3
    
    # Ensure Hermitian (SCS numerical noise can make it slightly non-Hermitian)
    M = (M + M.conj().T) / 2.0
    
    return M


d = 6
stop = False
# Initialize a random maximally mixed Choi matrix as a starting point
# J_fixed = np.eye(d**2) / d 
# rho_fixed = np.eye(d**2) / (d**2)
# J2 = J_fixed
# rho = rho_fixed
for j in range(10000):
    if stop: 
        break
    print(j)
    
    rho_fixed = generate_random_density_matrix(d**2)
    P_reshuffle, P_PT, P_TrB = build_linear_maps(d)
    J_fixed = generate_random_boundary_ppt_channel(d, P_PT, P_TrB)
    rho = rho_fixed
    J2 = J_fixed
    
    U = get_random_breuer_hall_U(d)
    n = 3
    #print('d=', d)
    for i in range(n):
        # Optimize T1 while keeping T2 fixed
        _, J1 = optimize_channel('T1', J2, rho, d, U) # Discard prob.value
        min_val = np.real(np.trace(rho @ calculate_M_matrix(J1, J2, U, d)))
        _, J2 = optimize_channel('T2', J1, rho, d, U) # Discard prob.value
        min_val = np.real(np.trace(rho @ calculate_M_matrix(J1, J2, U, d)))
        
        d2 = d**2
        K1 = (P_reshuffle @ J1.flatten()).reshape((d2, d2), order='C')
        K2 = (P_reshuffle @ J2.flatten()).reshape((d2, d2), order='C')
        K_composed = K1 @ K2
        J_composed = (P_reshuffle.T @ K_composed.flatten()).reshape((d2, d2), order='C')
        J_PT = (P_PT @ J_composed.flatten()).reshape((d2, d2), order='C')
        
        if i % 4 == 0 and i != 0:
            #print(i)
            U_new = optimize_U_scalable(rho, J_PT, d, get_breuer_hall_U(d))
            val_new = np.real(np.trace(rho @ calculate_M_matrix(J1, J2, U_new, d)))
            if val_new < min_val:
                U = U_new
        
        #print('old min val', min_val)
        rho, min_val = update_rho(calculate_M_matrix(J1, J2, U, d))
    
        if min_val < -1e-3:
            print('weeee ddidd ittt!!!!!!1')
            print(i, min_val)
            print(J1, J2, U, rho)
            stop = True
            break
    print(min_val)
\end{lstlisting}

